% answers_short.tex —— 2008 年·全国大纲卷Ⅱ·文科 · 简短答案速查
% 由 tools/make_short_answers.py 自 answers.tex 的 \gksolhead/\gkanswer 提取，
% 并与 2008/2008_数学真题_全国2文_答案要点.md 逐题比对后生成，勿手改（改源 MD 后重跑）。

\begin{gknotice}
\textbf{使用说明}：本节只给最终答案便于快速核对，完整推导见解析版。答案为编写者独立推导，\textbf{非官方标准答案}。
\end{gknotice}

\gkpart{选择题}
\begin{center}\begin{tabular}{ccccccccccccc}
\toprule
题号 & 1 & 2 & 3 & 4 & 5 & 6 & 7 & 8 & 9 & 10 & 11 & 12 \\
\midrule
答案 & C & B & D & C & D & D & A & B & A & B & B & C \\
\bottomrule
\end{tabular}\end{center}

\gkpart{填空题}
\begin{itemize}[leftmargin=2.6em]
  \item[\textbf{13.}] $2$
  \item[\textbf{14.}] $420$
  \item[\textbf{15.}] $2$
  \item[\textbf{16.}] 两组相对的侧面分别为平行四边形；两个底面都是平行四边形。
\end{itemize}

\gkpart{解答题}
\begin{itemize}[leftmargin=2.6em]
  \item[\textbf{17.}] (1) $\sin C=\dfrac{16}{65}$；(2) $\triangle ABC$ 的面积为 $\dfrac83$
  \item[\textbf{18.}] $S_{20}=200$ 或 $330$（等差数列公差分别为 $0$ 或 $1$，题面未排除公差为零）。
  \item[\textbf{19.}] (1) $\dfrac15$；(2) $\dfrac{13}{125}$
  \item[\textbf{20.}] (1) 见解析版；(2) 二面角 $A_1-DE-B$ 的余弦为 $\dfrac1{\sqrt{126}}$
  \item[\textbf{21.}] (1) $a=2$；(2) $a\leqslant\dfrac65$
  \item[\textbf{22.}] (1) $k=\dfrac23$ 或 $k=\dfrac38$；(2) 四边形 $AEBF$ 面积的最大值为 $2\sqrt2$
\end{itemize}
